% Pursuing Stacks (cotes 134-2 to 134-8): archive cote 134-5, site page 70 (PDF page 71), typescript page 257.
% typescript. Draft compared: ps4.tex:3522-3584.
% AI-transcribed by Claude Opus 5.5 (claude-opus-5-5, adaptive thinking, effort high), Message Batches API, batch ps2 (msgbatch_01NVwr1Yuo8qJf5rykeGBXPK), 2026-11-12; not yet checked by eye.
\scanpage{134-5/70}
\mspage{257}

\uline{tible}), and that moreover $M_c$ \uline{generates} $M_{as}$ ?

\deleted{The question just met with before was:
3) What extra condition (if any) is needed for a\unsure{\xout{n}} totally aspheric
asphericity structure, in order that any aspheric object \inserted{$I$} of $M$ \xout{is} be
"totally aspheric" i.e. the map $I \longrightarrow e_M$ is in $\unsure{W}_a$ ?}

Before looking up a little these questions, I would like \xout{first} \inserted{however} to
carry through at once the "idyllic picture" of canonical modelizers,
\deleted{as} forshadowed in section 50 (p.110), as I feel that this should be pos-%
sible at present at no costs. Taking into, account the reflections of
the later sections 57 and 59, we get the following set-up.

Let $M$ be a $\underline{U}$-category\inserted{stable under finite products}, endowed with a functor
(1) \quad $\hand{\pi_0^M \text{ or } \pi_0}$ \deleted{\illegible{1}} $: M \longrightarrow (\mathrm{Sets})$ ,

on which we make no assuption for the time being. We are thinking of
the example when $M$ is a totally 0-connected category (cf. prop. on page
142 for this notion) and $\hand{\pi}_o$ is the "connected components" functor, or \inserted{when}
$M = (\mathrm{Spaces})$ and $\hand{\pi}_o$ corresponds to taking sets of arc-wise connected
components. According to section 54 (p.1\over{\unsure{2}}{3}1), we introduce a correspon-%
ding homotopy interval structure h on $M$, admitting the generating family
of homotopy intervals made up with those intervals $\underline{I} = (I,\hand{\delta}_o,\hand{\delta}_1)$ \xout{such}
for which $I \longrightarrow e_M$ is "universally in $W_{\hand{\pi}_o}$" , namely
\[ \xout{a}\ \hand{\pi}_o(x \text{ x } I) \longrightarrow \hand{\pi}_o(x) \quad \text{bijective for any } x \text{ in } M \ . \]
(Under suitable conditions on $\hand{\pi}_o$ , this \xout{\unsure{structu}} homotopy interval
structure h on $M$ is the widest one "compatible with $\hand{\pi}_o$" in the sense
of page 130 , namely such that $\hand{\pi}_o$ transforms homotopisms into homoto-%
pisms - cf . proposition p. 133\unsure{\xout{7}}.) We are interested in the case when
this homotopy structure on $M$ can be described by a contractibility struc-%
ture $M_c$ on $M$ , which is then uni\unsure{qu}e of course, hence well defined in
terms of $\hand{\pi}_o$. Therefore the asphericity structure \xout{defined by} generated
by $M_c$ is equally well defined in terms of $\hand{\pi}_o$ , and likewise the corres-%
ponding notion $W_a$ of "weak equivalence" . We then get a canonical func-%
tor
\[ (2) \qquad \mathrm{Hot}_M = W_a^{-1}M \longrightarrow (\mathrm{Hot}) \]
(section 77, p.239). We'll have to find still a suitable \inserted{extra} condition on
the functor $\hand{\pi}_o$ , implying that this functor is canonically isomorphic
to the functor deduced from (2) by composing with the canonical functor

\margin{70}
